\fsection{Ej 2 Pág 136}
\fsubsection
[\textcolor{waveBlue2}{Tarea: Brecha de Seguridad}]
{\textcolor{waveBlue2}{\boldit{Tarea: Brecha de Seguridad}}}
{
	Elige un tipo de empresa relacionada con el sector profesional de tu ciclo formativo. Realiza
	un listado de los datos que gestiona, imagina una brecha de seguridad y define el procedimiento
	que debe llevarse a cabo con las personas afectadas y con la Agencia Española de Protección
	de Datos (AEPD), indicando: la categoría de los datos filtrados, la categoría y vulnerabilidad
	de las personas afectadas, el volumen de datos afectados y el procedimiento con la AEPD y los
	interesados.
}

\begin{solution}

	\textbf{Empresa:} \textit{AutoSistemas Industriales S.L.}

	\medskip
	\hrule
	\medskip

	\textbf{1. Tipos de datos que gestiona la empresa}

	\begin{itemize}
		\item Datos de empleados: nombre, DNI, número SS, salario, cuenta bancaria, partes de baja.
		\item Datos de clientes: nombre/razón social, NIF/CIF, dirección, teléfono, correo, IBAN.
		\item Datos de proveedores: nombre de contacto, teléfono, correo corporativo, IBAN.
		\item Datos de procesos internos: registros de intervenciones técnicas, accesos a
		      instalaciones de clientes.
		\item Listados de compras y ventas, presupuestos, facturas.
	\end{itemize}

	\medskip

	\textbf{2. Descripción de la brecha de seguridad (supuesto)}

	Un empleado del departamento administrativo recibe un correo de \textit{phishing} que simula ser
	una actualización del software de gestión. Al hacer clic en el enlace e introducir sus
	credenciales, un atacante externo accede al sistema durante 48 horas sin ser detectado.
	Durante ese tiempo, descarga la base de datos de clientes y el fichero de nóminas de empleados.
	El acceso es descubierto al detectarse un inicio de sesión desde una IP extranjera.

	Tipo de brecha: \textbf{brecha de confidencialidad} (acceso no autorizado a datos personales).

	\medskip

	\textbf{3. Categorías de datos afectados}

	\begin{itemize}
		\item Datos identificativos y económicos de \textbf{85 clientes} (nombre, NIF, IBAN,
		      dirección, correo).
		\item Datos laborales y económicos de \textbf{12 empleados} (DNI, número de afiliación SS,
		      salario, cuenta bancaria). Los datos de salud (partes de baja) no estaban en el fichero
		      comprometido.
	\end{itemize}

	\medskip

\end{solution}
\clearpage
\begin{solution}
	\textbf{4. Categoría de personas afectadas y su vulnerabilidad}

	\begin{itemize}
		\item \textbf{Clientes:} riesgo de usurpación de identidad, fraude bancario (datos IBAN
		      expuestos) y ataques de ingeniería social (suplantación de la empresa).
		\item \textbf{Empleados:} riesgo elevado por exposición de datos bancarios y número de
		      Seguridad Social, susceptibles de ser usados para fraude o suplantación laboral.
		      No se trata de colectivos especialmente vulnerables (menores, personas con discapacidad),
		      pero la naturaleza de los datos los sitúa en \textbf{riesgo alto}.
	\end{itemize}

	\medskip

	\textbf{5. Volumen de datos afectados}

	\begin{itemize}
		\item 97 personas físicas afectadas en total (85 clientes + 12 empleados).
		\item Estimación de registros: aprox. 970 campos de datos personales comprometidos.
	\end{itemize}

	\medskip

	\textbf{6. Procedimiento a seguir}

	\textbf{a) Medidas inmediatas de contención:}
	\begin{enumerate}[label=\arabic*.]
		\item Revocar inmediatamente las credenciales comprometidas y forzar cambio de contraseñas.
		\item Aislar los sistemas afectados y preservar los registros (logs) como evidencia.
		\item Documentar el incidente: hora de detección, sistemas afectados, acciones tomadas.
	\end{enumerate}

	\textbf{b) Notificación a la AEPD (art. 33 RGPD) -- plazo máximo: 72 horas:}
	\begin{enumerate}[label=\arabic*.]
		\item Acceder a la sede electrónica de la AEPD y cumplimentar el formulario de notificación
		      de brechas de datos personales.
		\item Indicar: naturaleza de la brecha, categorías y número de interesados afectados,
		      categorías de datos comprometidos, consecuencias probables, medidas adoptadas.
		\item La notificación puede realizarse de forma provisional si aún no se dispone de todos
		      los datos, completándola posteriormente.
	\end{enumerate}

	\textbf{c) Comunicación a los afectados (art. 34 RGPD) -- sin dilación indebida:}

	Dado que existe \textbf{alto riesgo} para los derechos y libertades (datos bancarios expuestos),
	es \textbf{obligatorio} comunicar la brecha a cada persona afectada. La comunicación debe incluir:
	\begin{itemize}
		\item Descripción clara de qué ha ocurrido y en qué período.
		\item Datos e información personal que se ha visto comprometida.
		\item Consecuencias posibles (fraude, suplantación de identidad).
		\item Medidas adoptadas por la empresa para mitigar el daño.
		\item Recomendaciones concretas para los afectados (vigilar movimientos bancarios,
		      cambiar contraseñas, activar alertas de fraude).
		\item Datos de contacto del responsable o del DPD para resolver dudas.
	\end{itemize}

	\textbf{d) Registro interno del incidente (art. 33.5 RGPD):}

	Con independencia de la obligación de notificar, la empresa debe documentar la brecha en su
	registro interno de incidentes, incluyendo hechos, efectos y medidas correctoras.

\end{solution}
